\section{Constructing the fixed-vector rule}\label{sec:algorithm}

We construct the rule of Proposition~\ref{prop:hidden} for a fixed hidden
weight vector. It uses geometric rounding and chooses with equal
probability between a single-choice maximum branch and a main branch.
The maximum branch handles weight concentrated on one label. The main
branch filters arriving labels against revealed weights and assigns
eligible arrivals to layers of nested flats constructed from those
revealed weights. Greedy independent sets from different layers can
then be combined feasibly.

\subsection{Rounding and a maximum branch}

Fix the constants
\begin{equation}\label{eq:constants}
 B=2^{32},\qquad \kappa=2^{100},\qquad t=2^{-140},\qquad
 \eta=2^{-12},\qquad \gamma=\eta/16=2^{-16}.
\end{equation}
Here $B$ separates weight levels, $\kappa$ is the rank charge in the
density expansions, and $t$ is the common sparsity of the guard and
eligibility masks. The exit-probability cutoff $\eta$ gives the
sample-to-rank coefficient $\gamma$. The choices leave room for the
losses combined in Section~\ref{sec:accounting}.
For $w_e>0$, round down to $u_e=B^{\lfloor\log_B w_e\rfloor}$, and put
$u_e=0$ when $w_e=0$. Define $u(Z)=\sum_{e\in Z}u_e$ and
$W=\OPT_M(u)$. The rule rounds each coordinate when its weight is
revealed or arrives; the full vector $u$ is used only for analysis. Then
\begin{equation}\label{eq:rounding}
 \OPT_M(w)/B\le W\le\OPT_M(w).
\end{equation}
Larger rounded weights have higher priority, with ties resolved by the fixed
label order. Priority greedy on all positive labels produces an independent
set $I^*$ of weight $W$: its intersection with every initial segment of
weight levels has full rank there, so summing the differences of successive
levels proves optimality.

Let $m_*$ be the largest rounded weight of a positive nonloop, or zero if
there is none. A \emph{maximum rule} reveals a fair independent mask and
accepts the first unmasked positive nonloop of higher priority than every
revealed positive nonloop, then accepts nothing further. If there are at
least two positive nonloops, with probability $1/4$ the highest-priority
one is unmasked and the second is masked. On this event only the highest
can exceed the threshold, so it is accepted in every order. With one
positive nonloop, success has probability $1/2$. Thus in all cases
\begin{equation}\label{eq:maximum}
 \Ex\big[\min_\pi u(A_{\mathrm{max}}(R,\pi))\big]\ge m_*/4.
\end{equation}
The final fixed-vector rule uses this rule or the main rule constructed
next, with equal probability. The maximum branch gives the guarantee when
the largest rounded weight is large enough relative to the rounded optimum;
the main branch handles the complementary case. Draw all coins, including unused branch
coins, initially and independently of the weights.

\subsection{Candidates and weight groups}

For the main rule draw mutually independent masks $H,D,C,T$, independently
across labels, with respective membership probabilities $1/2,1/4,t,t$.
Their roles are summarized below.
\begin{center}
\begin{tabular}{@{}ccl@{}}
\toprule
Mask & Probability & Role \\
\midrule
$H$ & $1/2$ & Filters candidates by higher-priority span \\
$D$ & $1/4$ & Lists weight groups and supplies density data \\
$C$ & $t$ & Supplies guards for the flat paths \\
$T$ & $t$ & Thins online eligibility independently of the paths \\
\bottomrule
\end{tabular}
\end{center}
Reveal the fixed-vector weights on $H\cup D\cup C$
and sacrifice these labels. Under Lemma~\ref{lem:simulation}, these revealed
weights are the corresponding sample coordinates.
Membership in $T$ alone reveals no weight and forces no rejection.
For each label $e$, let $H_{>e}$ be the labels in $H$ with higher priority
than $e$. Define the candidate set and a comparison independent set by
\begin{equation}\label{eq:candidates}
 U=\{e\notin H:u_e>0,\ e\notin\cl(H_{>e})\},
 \qquad Y=I^*\setminus H.
\end{equation}
The full set $U$ is used for analysis; the candidate test itself is
computable as soon as the weight of $e$ is known.

\begin{lemma}[Candidate mass]\label{lem:filter}
The set $Y$ is independent, $Y\subseteq U$, and
\begin{equation}\label{eq:candidate-mass}
 \Ex_H u(Y)=W/2,\qquad \Ex_H u(U)\le W.
\end{equation}
\end{lemma}
\begin{proof}
For $e\in I^*$, no higher-priority labels span $e$, so neither do the
higher-priority labels of $H$. This proves $Y\subseteq U$; its expectation
follows from the fair mask. Conditional on all higher-priority memberships
in $H$, the event $e\notin\cl(H_{>e})$ is fixed. Its contribution to $u(U)$
has the same expected weight as the event that $e$ belongs to $H$ and is
selected by priority greedy on $H$. Sum over the positive labels.
The resulting expected greedy weight on $H$ is at most $W$.
\end{proof}

Conditional on $H$, the \emph{true weight groups}
\[
 U_i=\{e\in U:u_e=B^i\},\qquad n_i=|U_i|,\qquad Y_i=Y\cap U_i
 \quad(i\in\mathbb Z)
\]
are fixed. Empty groups may be omitted. A group is \emph{listed} when
$D\cap U_i\ne\varnothing$. Enumerate the listed groups as $G_1,\ldots,G_m$
from lower to higher weight, and write $D_h=D\cap G_h$, $C_h=C\cap G_h$,
and $T_h=T\cap G_h$. The complete $D_h,C_h$ are known before arrivals:
every label of $D\cup C$ is revealed, and its candidate test uses only $H$
and its own weight. The algorithm needs neither the complete $G_h$ nor
the true group sizes $n_i$. All candidates are nonloops.

\subsection{Density expansions}

For a listed group $h$ and a flat $P$, let $Q_h(P)$ be the largest flat
maximizing
\begin{equation}\label{eq:density-objective}
 f_h(Z)=|D_h\cap Z|-\kappa\rk(Z)
 \qquad\text{over flats }Z\supseteq P.
\end{equation}
For a nonflat argument, take its closure first. The objective rewards the
number of $D_h$ labels spanned and charges $\kappa$ per unit of rank.
Comparison with a further extension bounds its additional sample count
by $\kappa$ times its additional rank. A separate telescoping argument in
Lemma~\ref{lem:generators} bounds the generators used over the entire path.

This is a contracted, sample-restricted form of the cardinality-minus-rank
maximization associated with principal partitions; see Santiago, Sergeev,
and Zenklusen~\cite[Section 3.1, Equation (1)]{SSZ25}. We prove the needed
monotonicity and residual-rank bounds directly.

\begin{lemma}[Density expansion]\label{lem:density}
The largest maximizer in \eqref{eq:density-objective} exists, the map $Q_h$
is increasing, and $P\subseteq Q_h(P)$. If $Q=Q_h(P)$, then for every set $Z$,
\begin{equation}\label{eq:density-residual}
 |D_h\cap(\cl(Q+Z)\setminus Q)|\le\kappa\rk(Z\mid Q).
\end{equation}
\end{lemma}
\begin{proof}
The function $f_h$ is supermodular on the lattice of flats: the counting
term is supermodular under intersection and closed union, and the negative
rank term is supermodular by \eqref{eq:rank-submodular}. There are finitely
many flats. The closed union of any two maximizers containing $P$ is again
a maximizer, so there is a largest one.

For $P\subseteq P'$, put $Q=Q_h(P)$ and $Q'=Q_h(P')$. The flat
$Q\cap Q'$ is feasible at $P$ and $\cl(Q\cup Q')$ is feasible at $P'$.
Optimality gives one inequality and supermodularity the reverse:
\[
 f_h(Q\cap Q')+f_h(\cl(Q\cup Q'))
       \le f_h(Q)+f_h(Q')
       \le f_h(Q\cap Q')+f_h(\cl(Q\cup Q')).
\]
Equality in the two optimality comparisons shows that $\cl(Q\cup Q')$
is also a maximizer at $P'$. Its largest maximizer contains $Q$, as claimed.
Finally compare $Q$ with the feasible flat $\cl(Q+Z)$ in
\eqref{eq:density-objective}. The difference of their ranks is
$\rk(Z\mid Q)$, giving \eqref{eq:density-residual}.
\end{proof}

\subsection{Nominal paths and guards}

Integer $k$ denotes an \emph{auxiliary time}, unrelated to arrival time.
Set $F_0(k)=L$ for $k<0$ and $F_0(k)=E$ for $k\ge0$. For each listed group,
put $a_h=-3h$ and define an enabled expansion
\[
 Q_{h,k}(P)=
 \begin{cases}Q_h(P),&k\ge a_h,\\ P,&k<a_h,\end{cases}
 \qquad(P\text{ a flat}).
\]
In increasing group order, form its \emph{nominal} and \emph{guarded} paths
\begin{equation}\label{eq:paths}
 X_h(k)=Q_{h,k}(F_{h-1}(k)),\qquad
 F_h(k)=\cl\bigl(X_h(k)+(C_h\cap X_h(k+1))\bigr).
\end{equation}
Thus a guard uses a label one auxiliary step before it enters the nominal
path. Each density map uses only $D_h$ and the known matroid, and each
guard set uses the revealed labels $C_h$. Consequently these flats can
be computed as subsets of $E$ before arrivals, even though the weights
of some labels they contain remain unknown.

\begin{lemma}[Path properties]\label{lem:paths}
The paths increase with $k$ and satisfy
\begin{equation}\label{eq:path-inclusions}
 F_{h-1}(k)\subseteq X_h(k)\subseteq F_h(k)\subseteq X_h(k+1).
\end{equation}
They equal $E$ for $k\ge0$. Moreover $X_h(k)=L$ for $k<a_h$ and
$F_h(k)=L$ for $k<a_h-1$. If $e$ is a nonloop, its nominal birth
\[
 b_h(e)=\min\{k:e\in X_h(k)\}
\]
is either $a_h$, with $e\in Q_h(L)$, or belongs to $[a_h+2,0]$.
\end{lemma}
\begin{proof}
Induct on $h$. The map $Q_{h,k}$ is increasing both in its argument and
in $k$, since enabling it replaces the identity by an extensive increasing
map. Thus $X_h$ and then $F_h$ are increasing. The first two inclusions
in \eqref{eq:path-inclusions} are extensiveness; the last follows because
both sets adjoined in the definition of $F_h(k)$ lie in the flat $X_h(k+1)$.
The same recursion gives $E$ at times at least zero.

For $h>1$, $a_{h-1}=a_h+3$ and the inductive lower bound makes
$F_{h-1}(k)=L$ for $k<a_h+2$; this also holds for $h=1$ directly from
$F_0$. Before $a_h$ the new map is the identity, proving the asserted
lower bounds for $X_h$ and $F_h$. At $a_h$ and $a_h+1$, the nominal flat
is the same flat $Q_h(L)$, so no nonloop is born at $a_h+1$.
\end{proof}

\subsection{The online decisions}

Choose an independent fair parity $\varepsilon\in\{0,1\}$ and put
$\widehat F(k)=F_m(k)$. When $m=0$, use $\widehat F=F_0$ and reject all
arrivals. In general the layer endpoints are the integers
$b\equiv\varepsilon\pmod2$, with layer $(b-2,b]$ having base
$\widehat F(b-2)$.
The rule maintains a set $A_b\subseteq\widehat F(b)$ in each layer,
independent over its base $\widehat F(b-2)$.
Since successive layers share an endpoint, all labels selected in an
earlier layer lie in the base of every later layer. Greedy selection in
alternating restriction--contraction minors also appears in
Feldman, Svensson, and Zenklusen~\cite[Section 3]{FSZ15}; here the nested
flats come from the guarded density paths.

At arrival of a label $e$, perform the following operations.
\begin{enumerate}
 \item Reject if $e\in H\cup D\cup C$, if its candidate test fails,
 or if its rounded weight is not listed.
 \item For its listed group $h$, compute $b=b_h(e)$. Reject unless
 \begin{equation}\label{eq:eligibility}
 e\in T,\qquad b\ge a_h+2,\qquad b\equiv\varepsilon\pmod2,
 \qquad e\notin\widehat F(b-2).
 \end{equation}
 \item Let $A_b$ be the previously accepted labels assigned to this
 layer. Accept $e$ exactly when
 $\rk(\{e\}\mid\widehat F(b-2)+A_b)=1$, and then add it to $A_b$.
\end{enumerate}
These operations use the revealed data, the arriving label and weight,
the seed, and the previous decisions. In particular, no full true weight
group or future arrival is needed.
By Lemma~\ref{lem:paths}, the three-step activation schedule leaves no
birth at $a_h+1$. Excluding the baseline birth $a_h$ therefore leaves exactly
the births for which $b-2\ge a_h$: the density map is already enabled
at the preceding layer endpoint. Figure~\ref{fig:paths} shows both the
guard step and the nesting that makes layerwise decisions feasible.

\begin{figure}[htbp]
\centering
\begin{tikzpicture}[>=Latex, every node/.style={font=\small},
 flat/.style={draw=linkblue,rounded corners=2pt,inner sep=7pt}]
 \node[flat] (prior) {$F_{h-1}(k)$};
 \node[flat,right=23mm of prior] (nominal) {$X_h(k)$};
 \node[flat,right=31mm of nominal] (guarded) {$F_h(k)$};
 \node[above=13mm of guarded] (guard) {$C_h\cap X_h(k+1)$};
 \draw[->] (prior) -- node[above] {$Q_{h,k}$} (nominal);
 \draw[->] (nominal) -- node[below] {adjoin and close} (guarded);
 \draw[->] (guard) -- (guarded);
 \node[flat,below=27mm of prior] (leftbase) {$\widehat F(b-2)$};
 \node[flat,below=27mm of nominal] (middlebase) {$\widehat F(b)$};
 \node[flat,below=27mm of guarded] (rightbase) {$\widehat F(b+2)$};
 \draw[->] (leftbase) -- node[above] {$\subseteq$}
     node[below=1mm,align=center] {layer $(b-2,b]$\\accept $A_b$} (middlebase);
 \draw[->] (middlebase) -- node[above] {$\subseteq$}
     node[below=1mm,align=center] {layer $(b,b+2]$\\accept $A_{b+2}$} (rightbase);
\end{tikzpicture}
\caption{Above: a density expansion followed by guards from the next
nominal time gives $F_h(k)\subseteq X_h(k+1)$. Below: two consecutive
layers of the chosen parity on the final path. The set $A_b$ is
independent over $\widehat F(b-2)$ and lies in $\widehat F(b)$;
$A_{b+2}$ is independent over this latter base. In the lower panel,
the horizontal order is auxiliary time, regardless of the actual
arrival order.}
\label{fig:paths}
\end{figure}

\begin{lemma}[Feasibility and measurability]\label{lem:feasibility}
The main rule maintains an independent accepted set at every arrival
prefix, for every permutation. It is a measurable rule based only on the
allowed information.
\end{lemma}
\begin{proof}
Every eligible label at endpoint $b$ lies in
$X_h(b)\subseteq\widehat F(b)$. The accepted set of its layer is
independent over $\widehat F(b-2)$. In increasing auxiliary endpoint order,
all earlier accepted layers lie in that base flat, so extending by the
current layer preserves independence. This proves independence of the
union for any interleaving of arrivals, and of every prefix.

All masks are drawn initially and the matroid is known. Each density
maximization is over finitely many flats. By Lemma~\ref{lem:paths}, it is
enough to compute a finite interval of integer times for the at most
$|E|$ listed groups. Positive rounding has countably many measurable
level sets; each realized vector has finitely many occupied levels.
All remaining operations are finite rank tests, comparisons, and set
operations with fixed tie rules. They therefore define measurable
decisions. Zeros are never passed to a logarithm, and loops fail the
candidate test.
\end{proof}

\subsection{A generator budget for the entire path}

The rule is now fully specified. To analyze it, we need a bound on how
much each density expansion can change the path. Bounding the rank
increase separately at each time would charge the same sample labels
repeatedly. The next lemma instead chooses one set of generators for all
times. It also bounds the independent sampled mass excluded by the
baseline-birth condition.

\begin{lemma}[A generator budget for the entire path]\label{lem:generators}
For each listed group $h$ there is a set $J_h\subseteq D_h$ and increasing
subsets $J_h(k)\subseteq J_h$ such that
\begin{equation}\label{eq:generator-budget}
 X_h(k)=\cl(F_{h-1}(k)+J_h(k)),\qquad
 |J_h|\le |D_h|/\kappa.
\end{equation}
In particular $\rk(Q_h(L))\le |D_h|/\kappa$.
\end{lemma}
\begin{proof}
Before $a_h$ no generators are needed. Process the enabled times
$k=a_h,a_h+1,\ldots,0$ in order. Given the preceding generators, set
\[
 P_k=\cl(F_{h-1}(k)+X_h(k-1)).
\]
By monotonicity, $P_k\subseteq X_h(k)$. It is feasible in the maximization
defining $X_h(k)$, whence
\begin{align*}
 \kappa\bigl(\rk(X_h(k))-\rk(P_k)\bigr)
 &\le |D_h\cap X_h(k)|-|D_h\cap P_k|\\
 &\le |D_h\cap X_h(k)|-|D_h\cap X_h(k-1)|.
\end{align*}
The flat $R_k=\cl(F_{h-1}(k)+(D_h\cap X_h(k)))$ is contained in
$X_h(k)$ and has exactly the same $D_h$ count. If this inclusion were
proper, $R_k$ would have smaller rank, since a proper subflat has
smaller rank. It would then improve the objective, contradicting
optimality. Thus $R_k=X_h(k)$, and the extension over $P_k$ has a
basis chosen from $D_h\cap X_h(k)$.

Adjoin such a basis to the generators, choosing by the fixed label order.
Since $F_{h-1}$ increases, the preceding generators together with
$F_{h-1}(k)$ already generate $P_k$. This proves the first identity
inductively. The new generators are distinct, and their counts satisfy
the displayed inequality. Summing it telescopes the $D_h$ counts and
gives $|J_h|\le |D_h|/\kappa$. Extend the generator sets constantly
for $k>0$. At $k=a_h$, the lower-group flat is $L$, yielding the last
assertion.
\end{proof}

For the analysis, we first work on a group's own nominal path. The next
section shows that earlier groups' guards protect a constant fraction of
its independent sampled mass against lower-weight competition. We then
transfer this sampled count to unsacrificed rank. Later groups' path
expansions will be paid for by a single rank budget in
Section~\ref{sec:accounting}.
